\section{Resumen - Abstract}

Este proyecto consiste en el desarrollo de un producto de software en forma de juego \textbf{MMO RPG Maker},
que permitirá a los usuarios crear o unirse a mundos explorables diseñados por otros jugadores.
El principal desafío radica en construir una infraestructura escalable y robusta, capaz de sostener la conexión simultánea de cientos de jugadores en una misma instancia,
procesar miles de transacciones en tiempo real y garantizar la estabilidad frente a un gran volumen de conexiones concurrentes,
manteniendo siempre una experiencia de juego fluida y confiable.
\vspace{0.25cm}

\textit{This project consists of developing a software product in the form of an \textbf{MMO RPG Maker} game,
which will allow users to create or join explorable worlds designed by other players.
The main challenge lies in building a scalable and robust infrastructure capable of supporting the simultaneous connection of hundreds of players within the same instance,
processing thousands of transactions in real time, and ensuring stability under a high volume of concurrent connections,
while always maintaining a smooth and reliable gameplay experience.}

\section{Palabras clave - Keywords}

Multijugador, online, crear contenido, concurrencia, MMO.
\vspace{0.25cm}

\textit{Multiplayer, online, content creation, concurrency, MMO.}

\section{Introducción}

El presente proyecto tiene como objetivo el desarrollo de un universo Massively Multiplayer Online (MMO)
que permita a cientos de usuarios conectarse simultáneamente en un mundo virtual.
La propuesta se centra en brindar a los jugadores la posibilidad de crear sus propios mundos y diseñar aventuras unicas dentro de ellos,
en los que otros usuarios puedan unirse, participar y compartir experiencias. Además, el sistema busca fomentar la creación de contenido por parte de la comunidad,
ofreciendo herramientas para que los usuarios no solo puedan confeccionar sus mundos y aventuras, sino también promocionarlos y difundirlos dentro de la plataforma.

\section{Estado del arte}

El desarrollo de juegos de rol multijugador en línea (MMO RPG) ha sido ampliamente explorado por títulos como World of Warcraft,
RuneScape y Final Fantasy XIV, los cuales sentaron las bases de los mundos persistentes con miles de jugadores conectados simultáneamente.
Sin embargo, estas propuestas se centran en infraestructuras cerradas que no permiten a los usuarios diseñar sus propios entornos,
limitando la personalización y la capacidad de experimentación dentro del ecosistema de juego. 
\vspace{0.25cm}

Por otro lado, motores como RPG Maker, Unity y Unreal Engine han democratizado la creación de videojuegos,
aunque su alcance técnico se orienta más hacia experiencias de un solo jugador o de multijugador reducido,
sin ofrecer de forma nativa soluciones que soportan un alto volumen de transacciones y conexiones en tiempo real.
\vspace{0.25cm}

En los últimos años, plataformas como Roblox y Minecraft han marcado un punto de inflexión al integrar la creación de mundos por parte de los jugadores
con experiencias multijugador de gran escala, convirtiéndose en referentes del contenido generado por usuarios (UGC). 
\vspace{0.25cm}

A nivel tecnológico, el avance de infraestructuras en la nube, como \textit{AWS GameLift}, \textit{Microsoft PlayFab} o \textit{Photon}, ha permitido que los desarrolladores 
gestionen de manera más eficiente la escalabilidad, el balanceo de carga y la persistencia de datos en tiempo real.
\vspace{0.25cm}

No obstante, aún existe un espacio de innovación en la intersección entre la accesibilidad de herramientas de creación como RPG Maker 
y la robustez de arquitecturas masivas propias de los MMO RPG, lo cual abre la oportunidad para proyectos que busquen combinar ambas dimensiones en una misma propuesta.

\section{Problema detectado y/o faltante}

\subsection{Estilo del Producto}

Un aspecto fundamental para diferenciarnos de otras plataformas de creación de contenido, como por ejemplo, Minecraft o Roblox, es ofrecer un estilo visual simple,
coherente y fácil de modificar. La simplicidad no solo debe reflejarse en la apariencia, sino también en las herramientas de personalización,
de manera que cualquier usuario pueda personalizar los artefactos de su mundo sin barreras técnicas.

\subsection{Creación de mundos por usuarios}
Uno de los pilares del proyecto es ofrecer a los jugadores la posibilidad de crear y compartir sus propios mundos explorables.
Sin embargo, este enfoque plantea una serie de desafíos:

\begin{itemize}
    \item \textbf{Accesibilidad de la interfaz:} ¿Cómo diseñar un editor que sea intuitivo y fácil de usar para personas sin conocimientos de programación
    o de diseño de videojuegos, pero que al mismo tiempo ofrezca profundidad para usuarios avanzados?
    \item \textbf{Equilibrio entre abstracción y personalización:} ¿Cómo proporcionar herramientas de creación simplificadas que permitan generar contenido rápidamente,
    sin sacrificar la capacidad de personalización y creatividad individual?
    \item \textbf{Limitaciones del editor:} ¿Qué restricciones tendrá el creador de mapas? ¿Será posible recrear experiencias comparables a MMOs
    reconocidos con las herramientas disponibles, o se priorizará un alcance más reducido y específico?
\end{itemize}

\subsection{Incentivo de Jugar y Crear}
Otro aspecto central del proyecto es definir qué motiva tanto a los creadores de mundos como a los jugadores a participar activamente en la plataforma.
Es necesario establecer mecanismos de incentivo que fomenten la creación de contenido de calidad y la interacción constante dentro de los mundos generados. 
\vspace{0.25cm}

Entre las preguntas clave a resolver se encuentran: ¿habrá algún tipo de “premio” o recompensa por jugar?
¿Los creadores podrán obtener un beneficio económico por sus mundos? 
¿Qué elementos harán que los jugadores quieran permanecer en un mundo específico, garantizando que este sea atractivo, disfrutable y mantenga su calidad a lo largo del tiempo?

\subsection{Mantenimiento y Monetización}
Otro de los principales problemas es la monetización y mantenimiento a largo plazo del proyecto. La monetización es el mecanismo de sustento económico,
por el cual se podrá mantener durante mucho tiempo funcionando el proyecto, es la base de sustentación operativa elemental del proyecto.
\vspace{0.25cm}

La cuestión a enfrentarse implica como cobrar dentro del proyecto, como así también qué porcentajes.
Cómo impacta cada una de las decisiones de cobro, cómo el contexto afecta a cada una de las políticas, cómo reaccionan cada uno de los agentes ante cambios en las políticas de cobro.
\vspace{0.25cm}

También nos enfrentamos a cuestiones relacionadas con los incentivos dentro de la plataforma.
Los incentivos y sus cambios afectan al uso a lo largo del tiempo de la plataforma, lo que impacta en la durabilidad del MVP a largo plazo.
Un mal cambio en las políticas de incentivo y cobros, puede llevar a la pérdida de interés en la plataforma y llevar al fin del proyecto.
\vspace{0.25cm}

Por eso, se hace relevante la implementación de buenas políticas de incentivos y cobros, que no impacten de forma negativa en los costos del proyecto,
los incentivos para crear contenido dentro de la plataforma y la jugabilidad y uso del mismo por parte de los usuarios.

\subsection{Arquitectura e Infraestructura}
Uno de los principales desafíos del proyecto es definir una arquitectura capaz de soportar la complejidad de un sistema multijugador masivo 
con mundos generados por los usuarios. A diferencia de un MMO RPG tradicional —donde existe un único mundo persistente gestionado por servidores dedicados de tipo zoning—,
esta propuesta requiere sostener múltiples mundos de manera simultánea, cada uno con baja latencia y alto rendimiento para garantizar una experiencia de juego fluida.
\vspace{0.25cm}

Además, se deben considerar aspectos clave como:

\begin{itemize}
    \item \textbf{Escalabilidad:} ¿Cómo gestionar de manera eficiente cientos de clientes conectándose a un mismo mundo, y miles distribuidos en diferentes instancias?
    \item \textbf{Baja latencia:} Garantizar tiempos de respuesta mínimos en un entorno de interacción en tiempo real.
    \item \textbf{Costos de infraestructura:} ¿Qué inversión implica desplegar y mantener un sistema de estas características, considerando servidores,
    balanceadores, bases de datos distribuidas y servicios en la nube?
    \item \textbf{Orquestación de instancias de mundos:} Se requiere un sistema dinámico de autoescalado que levante o apague instancias según la demanda para evitar desperdicio de recursos.
    \item \textbf{Persistencia de datos en tiempo real:} Es necesario guardar progresos e inventarios en bases de datos distribuidas equilibrando consistencia y escalabilidad.
    \item \textbf{Tolerancia a fallos y continuidad del servicio:} La arquitectura debe ser resiliente, con replicación y recuperación automática ante caídas de servidores.
    \item \textbf{Seguridad y prevención de abusos:} El sistema debe mitigar ataques de \textit{cheating} mediante validación de tráfico y monitoreo continuo.
    \item \textbf{Costos operativos:} Los principales gastos estarán en servidores de baja latencia y almacenamiento persistente, requiriendo una planificación eficiente.
    \item \textbf{Estrategias de monetización ligadas a la infraestructura:} El modelo elegido (suscripción, micropagos o marketplace) determinará la asignación de recursos 
    por jugador y creador.
\end{itemize}

\subsection{Alcance del proyecto}
Este proyecto es, sin duda, ambicioso, y resulta evidente que no será posible resolver todos los desafíos en un solo año.
Sin embargo, sí es factible definir un MVP que cubra lo mínimo indispensable para abordar los problemas identificados previamente y que,
al mismo tiempo, siente las bases para optimizar y expandir el producto en futuras iteraciones.

\section{Solución propuesta}

\subsection{Estilo del Producto}
El estilo visual será unificado en toda la plataforma bajo un formato 2.5D. Los objetos de fondo y del mapa estarán representados en 3D,
mientras que los personajes y elementos animados se desarrollarán en 2D. Esta elección facilita tanto el diseño por parte del equipo como la creación de cosmeticos por los usuarios,
ya que un cosmético puede consistir simplemente en una imagen aplicada sobre el personaje o en animaciones 2D básicas (idle, movimiento, ataque).
De esta manera, se reduce la complejidad de implementar modelos 3D completamente animados, manteniendo al mismo tiempo un estilo visual atractivo 
y accesible para los usuarios.

\subsection{Creación de mundos por usuarios}
El creador de mundos será una aplicación que los usuarios podrán instalar para diseñar y personalizar mapas.
Al iniciar un nuevo proyecto, se definirá el tamaño del mapa (pequeño, mediano o grande) y, opcionalmente, se podrá importar un \textit{template} preexistente como punto de partida.
\vspace{0.25cm}

Dentro de cada mundo existirán tres tipos principales de entidades:

\begin{itemize}
    \item \underline{Assets}: objetos posicionables directamente en el mapa.
    \item \underline{Items:} elementos que los jugadores pueden almacenar en su inventario.
    \item \underline{Cosméticos:} objetos destinados a la personalización del personaje del jugador.
\end{itemize}

Cada creador podrá importar sus propias entidades para su mundo y brindarle una identidad única o, en su defecto,
utilizar las librerías de entidades pre-cargadas en la herramienta. Los mundos tendrán una estructura plana y cuadrada en entornos 3D,
con dimensiones proporcionales al tamaño de un jugador. A modo de ejemplo, si un jugador mide una unidad, un mapa considerado “pequeño” será de 500 unidades de extensión,
escalando dicho valor para los tamaños superiores.
\vspace{0.25cm}

Los usuarios podrán importar modelos 3D o 2D y asignarles comportamientos mediante interfaces abstractas. Por ejemplo, un modelo de criatura podrá configurarse como “Enemigo”,
definiendo su tipo de ataque o interacción con los jugadores. Este mismo esquema será aplicable a items y cosméticos, permitiendo ampliar la lógica de personalización y funcionalidad.
\vspace{0.25cm}

Asimismo, los creadores tendrán la posibilidad de diseñar \textbf{quests} o aventuras, que incentiven la exploración y la interacción con el mundo. De forma abstracta,
un \textit{quest} podrá dispararse mediante un disparador asociado a una acción, como recoger un objeto o interactuar con un NPC \textit{(non-playable-character)}. Al completarlo,
los jugadores recibirán recompensas definidas por el creador, tales como ítems o monedas.
\vspace{0.25cm}

En última instancia, cada creador será responsable de diseñar la propuesta de valor de su mundo, mientras que la plataforma proveerá las herramientas necesarias
para facilitar y potenciar dicha experiencia de creación.
\vspace{0.25cm}

Por último, se implementarán tiendas virtuales personalizables para cada mundo donde los jugadores podrán gastar sus monedas y/o gemas. 
Esto se detalla en mayor profundidad en la sección de \textbf{Mantenimiento y Monetización}.

\subsection{Mantenimiento y Monetización}
El sistema de monetización se basará en dos tipos de monedas virtuales, complementadas con un esquema de suscripciones para los creadores de mundos.

\begin{itemize}
    \item \textbf{Gemas:} moneda premium adquirida con dinero real, utilizable en cualquier mundo para obtener ítems y cosméticos por medio de la tienda.
    Este sistema está orientado a incentivar a los creadores a diseñar y ofrecer contenido estéticamente atractivo.
    \item \textbf{Oro:} moneda interna de cada mundo, no transferible entre mundos. Su generación es definida por los creadores a través de recompensas por quests,
    combates o eventos, y se utiliza para adquirir ítems de la tienda del mundo.
\end{itemize}

Cada creador podrá establecer tiendas internas en su mundo, definiendo precios en oro y/o gemas, así como el stock disponible de cada ítem.
Además, se implementará un sistema de suscripciones para creadores de mundos, que será necesario para mantener el mundo funcionando online.
Inicialmente, se ofrecerá un período de prueba gratuito de dos semanas \textit{(periodo tentativo)}, y posteriormente,
el costo de mantenimiento dependerá del tamaño y la demanda del mundo.
\vspace{0.25cm}

El objetivo es que los mundos estén bien diseñados, que incentiven la interacción y el consumo de ítems, de manera que puedan sostener sus propios costos gracias al gasto de los jugadores,
reduciendo así la carga financiera directa del creador.
\vspace{0.25cm}

Finalmente, los creadores tendrán la posibilidad de convertir gemas en dinero real, regulado por umbrales que deben superar para obtener dicha conversión,
estableciendo un esquema de ingresos directos. La plataforma retendrá un porcentaje de cada conversión como comisión por el servicio.

\subsection{Arquitectura e Infraestructura}
Se diseñó una arquitectura tentativa para los servidores que se consideran necesarios y sus comunicaciones.

Por un lado, deberá haber servidores escalables por zonas (de mundos), con el objetivo de soportar una gran cantidad de jugadores conectados en un mundo, así como también
deberá haber conjuntos de servidores (clusters), simbolizando los múltiples mundos simultáneos.
\vspace{0.25cm}

Por otra parte, un modelo de servidor con múltiples funcionalidades de tipo monolítico, donde estarán presentes distintos sistemas, como el sistema de cobranzas,
el manejo de datos de los mundos, el sistema del portal principal de mundos, etc.

Esto se desarrollará bajo ese tipo de arquitectura simplificada debido a que la mayor dificultad del proyecto se encuentra en la otra parte mencionada 
que llevará una arquitectura más compleja y que deberá optimizarse mucho para su correcto funcionamiento. De esta forma evitar sobre-complicar
partes del sistema que no generan valor al usuario, y que no son la característica principal del desafío planteado.

Se consideró al sistema elegido suficiente para el alcance del proyecto, y se espera crear un sistema bien elaborado y escalable.
\vspace{0.25cm}

Los usuarios cuentan con dos ejecutables principales: el cliente de juego, que les permite navegar y conectarse a distintos mundos, y el World Maker,
una aplicación destinada a los creadores para diseñar y publicar mundos. 

Ambos se comunican con la plataforma a través de un \textit{Gateway} central que actúa como punto de entrada único, manejando la seguridad y el ruteo de todas las peticiones.
\vspace{0.25cm}

Detrás del Gateway, la comunicación interna entre servicios se gestiona mediante un Message Broker. Este componente es el núcleo de la infraestructura,
ya que permite que cada servidor se conecte de forma desacoplada, habilitando así la escalabilidad independiente de cada parte del sistema.
Gracias a este diseño, si aumenta la cantidad de jugadores, se pueden lanzar más instancias de los servicios necesarios sin afectar al resto.
\vspace{0.25cm}

La simulación de los mundos corre en servidores de zona, que son instancias de Unity en modo headless, 
optimizadas únicamente para procesar lógica y sincronización en tiempo real. Estos servidores están clasificados en tres tipos según el tamaño de los mapas:
Clase A para mundos grandes, Clase B para mundos medianos y Clase C para mundos pequeños. Este esquema permite ajustar la capacidad de cómputo según la complejidad
y el volumen de jugadores de cada mundo, a la vez que facilita la escalabilidad horizontal, ya que se pueden lanzar múltiples instancias de cada clase según la demanda.
\vspace{0.25cm}

En conjunto, esta infraestructura busca equilibrar flexibilidad, rendimiento y costos. La nube y tecnologías como Kubernetes permiten orquestar los servidores 
de zona de manera dinámica, levantando y apagando recursos en función del tráfico. El uso de un message broker asegura que la comunicación entre servicios sea confiable 
y tolerante a fallos. 
El resultado es una base sólida para sostener un ecosistema de mundos creados por los usuarios, en donde la calidad de la experiencia depende tanto de la creatividad 
de los creadores como de la solidez de la arquitectura que lo soporta.

\subsection{Alcance del proyecto}

\section{Evaluación preliminar de impacto social y ambiental}

Se espera que el impacto social pueda generar una comunidad que se dedique a crear contenido innovador, que bajo otras circunstancias no podría darse,
ya que hoy en día crear nuevo contenido de videojuegos requiere la necesidad tiempo y dedicación para aprender a utilizar y familiarizarse con las distintas herramientas.
\vspace{0.25cm}

Por lo tanto, se busca que los usuarios puedan encontrarle rédito de distintas maneras, por un lado publicando el contenido en redes sociales, 
tanto como ofreciendo su producto, como así también de cómo lo hizo y las posibilidades que puedan tener otros usuarios.
Por otro lado, los usuarios de igual manera pueden consumir el contenido existente y publicar sus avances mediante servicios de streaming.
\vspace{0.25cm}

No hay una certeza absoluta en la que este proyecto, si una vez finalizado más allá del MVP que se entregará como proyecto final,
pueda a escalar más de lo esperado y se necesite continuo desarrollo de nuevo contenido bajo la demanda de la comunidad que pueda haber.
\vspace{0.25cm}

Además, todo esto depende de los usuarios creadores que utilicen el editor de nuestro proyecto para generar contenido atractivo para los usuarios,
por lo tanto debemos estar al tanto de las necesidades de esos creadores y estar solucionando todos los problemas presentes en el feedback
y demandas que se puedan disponer los usuarios.

\section{Metodología}

Para llevar a cabo este proyecto, se realizará una planificación bien estructurada sobre la gestión del mismo, previamente a comenzar el desarrollo,
para poder establecer primeramente su alcance, las problemáticas que puedan llegar a surgir, enumerar las investigaciones previas y las que se llevarán durante
el transcurso del desarrollo, y establecer puntos de entrega por iteraciones dentro de un tiempo establecido para llevar una meta concisa de los aspectos 
principales que debe poder satisfacer el proyecto en sí, con respecto a los objetivos planteados.
\vspace{0.25cm}

Esta planificación será basada en lo aprendido por las materias cursadas en la carrera de la facultad, utilizando las mejores herramientas necesarias para 
formar una buena gestión del proyecto que se llevará a cabo, y además para poder llevar una comunicación estable tanto entre los alumnos como con el tutor designado.
Por lo tanto, una vez establecido los objetivos principales, las características a desarrollar del proyecto y las metodologías a implementar para el desarrollo,
se dividirá estimativamente las tareas a realizar a cada uno de los integrantes del grupo, con el fin de avanzar en el trabajo lo más organizado posible.
\vspace{0.25cm}

Utilizaremos herramientas particulares como Miro, para plantear ideas generales de forma más visual e interactiva, como por ejemplo el desarrollo de User Story Maps
o planteos basados en imágenes y diseños en croquis. Se ha creado hasta el momento en que se escribe este documento un user story map,
con detalles en usuarios basados como un creador de contenido y como de un usuarios jugador que consume contenidos, 
con un par de iteraciones con el fin de ordenar y definir ideas generales de las características que poseerá nuestro proyecto y sus funcionalidades más importantes.
\vspace{0.25cm}

Otra herramienta útil es GitHub Projects, que sirve para organizar las características por temas y plantear estimaciones semanales o bisemanales.
Esto permite establecer un cronograma claro para monitorear los avances del proyecto.
\vspace{0.25cm}

Aunque se crearán issues para detallar problemas por resolver —ya sean funcionalidades específicas, aspectos de arquitectura o tareas de investigación—,
estos estarán sujetos a modificaciones constantes. El objetivo es adaptarse tanto al alcance del proyecto como a la necesidad de coordinar las funcionalidades,
de modo que las dependencias entre ellas no retrasen el desarrollo.
\vspace{0.25cm}

En cuanto a riesgos del proyecto, principalmente se tiene en cuenta que el desafío está en brindar el servicio propuesto como se ha explicado
en el inciso de \textbf{Solución Propuesta}, donde se desarrollará un ambiente de creación \textit{amigable} a usuarios, pero esto es una estimación que para el proyecto final,
terminará siendo subjetivo, ya que hay una proporcionalidad entre la cantidad de características que se pueden brindar al usuario con la complejidad
que puede presentar ese editor mencionado, por lo tanto se intentará mantener la premisa de que sea un entorno de fácil uso y entendimiento,
pero que sin tener un desarrollo avanzado esto puede llegar a perder algo de valor en el producto final.
\vspace{0.25cm}

Otro de los riesgos presentes en este proyecto, es que es un desafío sobre la temática relacionada con una estructura de servidores que pueda ser estable
y capaz de soportar múltiples usuarios, teniendo en cuenta que no sólo depende de la calidad de desarrollo de dicho sistema,
sino que también hay factores de recursos que debemos disponer a la hora de poner a prueba,
por lo tanto la inversión en estos recursos es un factor a tener en cuenta para la prueba de rendimiento y estrés de nuestro sistema.
\vspace{0.25cm}

Un riesgo además es que el sistema monetario que presentemos, tiene que ser robusto, no sólo para que nosotros como desarrolladores salgamos beneficiados
y no tengamos pérdidas monetarias, sino para que los usuarios tampoco salgan perjudicados por otros usuarios con malas intenciones.
Por lo tanto hay que realizar un análisis exhaustivo y simulaciones para evitar posibles problemas.

\section{Experimentación y/o validación}

Este apartado describe el conjunto de pruebas validadas que se ejecutarán para verificar el cumplimiento de los objetivos técnicos y funcionales del proyecto. El plan de validación está diseñado para asegurar la robustez, escalabilidad y usabilidad del producto final, y estará sujeto a ajustes iterativos durante el desarrollo para adaptarse a la evolución del proyecto.


\subsection{Pruebas de Rendimiento y Estrés (Load Testing)}

El principal desafío técnico del proyecto, como se establece en la introducción, es la construcción de una infraestructura escalable y robusta.
Para validar este aspecto de manera cuantitativa, se ejecutarán conjuntos de pruebas de rendimiento y estrés centradas en los servidores de zona (Clase A, B y C).
El objetivo es medir la capacidad del sistema para sostener la conexión simultánea de cientos de jugadores en una misma instancia,
procesar miles de transacciones en tiempo real y garantizar la estabilidad bajo un alto volumen de conexiones concurrentes.
Además nos servirá para determinar la carga máxima al momento que se haga la prueba en la etapa final.

\begin{itemize}
    \item \textbf{Metodología:} Se utilizarán herramientas de profiling y simuladores de carga (como k6, Locust o JMeter) para emular el comportamiento de usuarios concurrentes.
    Las pruebas se realizarán de forma gradual, incrementando la carga desde decenas hasta cientos de usuarios virtuales por instancia de mundo.

    \item \textbf{Métricas Clave:} Se monitorizarán y documentarán rigurosamente los siguientes indicadores:
    \begin{itemize}
        \item \textbf{Latencia:} Tiempo de respuesta del servidor para acciones clave (movimiento, interacciones, chat).
        \item \textbf{Uso de CPU y Memoria:} Consumo de recursos de las instancias de los servidores de zona.
        \item \textbf{Tasa de Errores:} Porcentaje de requests fallidos o desconexiones bajo carga.
        \item \textbf{Escalabilidad:} Capacidad del sistema de orquestación para escalar automáticamente los servidores de zona según la demanda.
    \end{itemize}

    \item \textbf{Enfoque Iterativo:} Si bien las pruebas masivas se planifican para cuando se tenga un desarrollo sostenible del núcleo de networking,
    se implementará un protocolo de pruebas de carga progresivas (progressive load testing) desde las fases iniciales del desarrollo de esta feature.
    Esto permitirá detectar cuellos de botella y problemas de concurrencia de forma temprana.
\end{itemize}

\subsection{Pruebas de Usabilidad y Control de Calidad con Usuarios Externos (Alpha/Beta Testing)}

Dado que el proyecto tiene un fuerte componente comunitario y de creación de contenido, es crucial validar la usabilidad de las herramientas y la coherencia 
de la experiencia de usuario. Para ello, se reclutará a un grupo de usuarios externos al equipo de desarrollo, representativos del target de creadores y jugadores.

\begin{itemize}
    \item \textbf{Fase Alpha (Interna):} En una etapa temprana, se probarán las funcionalidades básicas del \textbf{World Maker} 
    (creación de mapas, importación de assets, configuración de entidades) con un grupo reducido y controlado de usuarios.
    El foco estará en la detección de bugs críticos y en la validación de los flujos de trabajo fundamentales.
    \item \textbf{Fase Beta (Pública o Cerrada):} En una etapa más avanzada, se ampliará el acceso a un grupo más grande de usuarios para probar el contenido generado.
    Los participantes crearán sus propios mundos, quests y tiendas, y otros jugadores probadores se encargarán de explorarlos. Se intentará obtener interesados 
    conocidos para probar todas estas características.
    \item \textbf{Objetivo:} Esta fase busca identificar inconsistencias, bugs específicos de interacción y problemas de diseño que no son evidentes para el equipo de desarrollo.
    El feedback se recopilará mediante formularios estructurados, sesiones de observación y herramientas de analytics integradas, 
    permitiendo ajustar la interfaz y la lógica para garantizar una experiencia intuitiva y libre de errores.
\end{itemize}


\subsection{Pruebas de Integración y Comunicación entre Servicios}

Para validar la solidez de la arquitectura propuesta, especialmente la comunicación entre el \textbf{Gateway}, el \textbf{Message Broker} y los distintos microservicios y servidores de zona,
se implementarán pruebas de integración automatizadas.

\begin{itemize}
    \item \textbf{Alcance:} Se simularán escenarios de fallo (ej.: caída de un servidor de zona, saturación del broker) para verificar la resiliencia del sistema 
    y su capacidad de recuperación.
    \item \textbf{Validación de Transacciones:} Se testeará exhaustivamente el flujo de transacciones en tiempo real, 
    como las compras en las \textbf{tiendas virtuales} de cada mundo, para garantizar la consistencia de datos entre el servidor de zona y los servicios monolíticos 
    de backend (ej: sistema de cobranzas).
\end{itemize}

\subsection{Suite de Pruebas de Regresión Automatizadas}

Se desarrollará una suite automatizada de pruebas end-to-end y de unit testing que se ejecutará de forma continua con cada integración de nuevo código (CI/CD).
Esto asegurará que los avances en el desarrollo y las nuevas funcionalidades no introduzcan errores en el código ya existente,
manteniendo la estabilidad del producto throughout todo el ciclo de desarrollo.
\vspace{0.25cm}

Este plan de validación integral asegurará que el producto final no solo sea funcional, sino también escalable, robusto y alineado con las expectativas de su comunidad de usuarios.

\section{Plan de actividades}

\section{Referencias}

\section{Anexo}
